\section{The Leontief input--output model}
The production vector $x$ must supply both intermediate consumption $Cx$
and final demand $d$. Thus
\[x=Cx+d,\qquad (I-C)x=d.\]
Column $j$ of $C$ lists the amounts each sector supplies as inputs per
unit produced in sector $j$. These columns need not sum to one; the
remaining part includes value added. This is a different model from the
closed exchange economy of Section 1.6.

The vector $d$ is the external demand, not the total required production.
When $I-C$ is invertible, $x=(I-C)^{-1}d$. For a single right side,
solving $(I-C)x=d$ by elimination is often more efficient than computing
an inverse. Respect the instructions against forming inverses in
Exercises 2 and 3 below.

\practice{2.6.P1}{Write each unit-consumption vector as a column}
The Goods column is $(.2,.5)^T$ and the Services column is $(.4,.3)^T$.
Hence
\[x=\begin{bmatrix}.2&.4\\.5&.3\end{bmatrix}x+
\begin{bmatrix}20\\30\end{bmatrix},\quad\text{or}\quad
\begin{bmatrix}.8&-.4\\-.5&.7\end{bmatrix}x=
\begin{bmatrix}20\\30\end{bmatrix}.\]
The first row accounts for all Goods required, and the second accounts
for all Services required. This sets up the requested model.

\exercise{2.6.1}{Construct the three-sector consumption matrix}
Use Manufacturing, Agriculture, Services as the order of rows and
columns. The input columns described in the exercise give
\[C=\begin{bmatrix}.10&.60&.60\\.30&.20&0\\.30&.10&.10\end{bmatrix}.\]
For example producing one agricultural unit requires $.60$ unit
from Manufacturing, $.20$ from Agriculture, and $.10$ from Services.
If Agriculture produces $100$ units, its intermediate-demand vector
is $100$ times that column:
\[(60,20,10)^T.\]
These are the inputs needed for that production plan, not a final-demand
vector or a complete equilibrium production plan for all sectors.

\exercise{2.6.2}{Meet agricultural final demand without computing an inverse}
The final-demand vector is $(0,20,0)^T$. Multiplying the equations
$(I-C)x=d$ by $10$ gives
\[\left[\begin{array}{rrr|r}9&-6&-6&0\\-3&8&0&200\\-3&-1&9&0\end{array}\right].\]
Swap rows 1 and 2. Then use $R_2\leftarrow R_2+3R_1$,
$R_3\leftarrow R_3-R_1$, and $R_3\leftarrow2R_3+R_2$.
The result is
\[\left[\begin{array}{rrr|r}-3&8&0&200\\0&18&-6&600\\0&0&12&200\end{array}\right].\]
Back substitution gives $x_S=50/3$,
$x_A=(600+6(50/3))/18=350/9$, and
$x_M=(8(350/9)-200)/3=1000/27$.
\answer{$x=(1000/27,350/9,50/3)^T$. Substitution into the original
model gives final demand $(0,20,0)^T$ exactly.}

\exercise{2.6.3}{Meet manufacturing final demand by the same elimination}
Now $d=(20,0,0)^T$, so the initial augmented right side after
multiplying by $10$ is $(200,0,0)^T$.
Use the same row swap and replacements listed in Exercise 2.6.2.
They produce
\[\left[\begin{array}{rrr|r}-3&8&0&0\\0&18&-6&200\\0&0&12&200\end{array}\right].\]
Thus $x_S=50/3$, $x_A=(200+100)/18=50/3$, and
$x_M=8(50/3)/3=400/9$.
\answer{$x=(400/9,50/3,50/3)^T$. No inverse matrix was used.}

\exercise{2.6.4}{Add the production responses to two demand vectors}
The new demand $(20,20,0)^T$ is the sum of the demands in the two
previous exercises. If $(I-C)x^{(a)}=d^{(a)}$ and
$(I-C)x^{(b)}=d^{(b)}$, additivity gives
$(I-C)(x^{(a)}+x^{(b)})=d^{(a)}+d^{(b)}$.
Add the actual production vectors:
\[x=\begin{bmatrix}1000/27\\350/9\\50/3\end{bmatrix}
+\begin{bmatrix}400/9\\50/3\\50/3\end{bmatrix}
=\begin{bmatrix}2200/27\\500/9\\100/3\end{bmatrix}.\]
The elimination pivots in Exercises 2 and 3 are nonzero, so the
coefficient matrix is invertible and this solution is unique.

\exercise{2.6.5}{Use the inverse requested in the problem}
Here
\[I-C=\begin{bmatrix}1&-1/2\\-3/5&4/5\end{bmatrix},\qquad
\det(I-C)=\frac45-\frac3{10}=\frac12.\]
The $2\times2$ inverse formula therefore gives
\[(I-C)^{-1}=\begin{bmatrix}8/5&1\\6/5&2\end{bmatrix}.\]
Multiplication by $(50,30)^T$ yields
\[x=(80+30,60+60)^T=(110,120)^T.\]
Check intermediate demand: $Cx=(60,90)^T$, so
$Cx+d=(60+50,90+30)^T=x$.

\exercise{2.6.6}{An inverse with a different determinant}
Now
\[I-C=\begin{bmatrix}4/5&-1/2\\-3/5&9/10\end{bmatrix},\qquad
\det(I-C)=\frac{36}{50}-\frac3{10}=\frac{21}{50}.\]
Consequently
\[(I-C)^{-1}=\frac{50}{21}
\begin{bmatrix}9/10&1/2\\3/5&4/5\end{bmatrix}
=\begin{bmatrix}15/7&25/21\\10/7&40/21\end{bmatrix}.\]
With $d=(16,12)^T$, the two coordinates are
$240/7+100/7=340/7$ and $160/7+160/7=320/7$.
\answer{$x=(340/7,320/7)^T$. This vector satisfies $(I-C)x=(16,12)^T$.}

\exercise{2.6.7}{The effect of increasing final demand by one unit}
Use the inverse from Exercise 2.6.5.
\begin{enumerate}[label=(\alph*)]
\item For demand $(1,0)^T$, the production vector is the first
column of the inverse: $(8/5,6/5)^T$.
\item For demand $(51,30)^T$, multiplication gives
$x=((8/5)51+30,(6/5)51+60)^T=(558/5,606/5)^T$.
\item The new demand is $(50,30)^T+(1,0)^T$. By distributivity of
matrix multiplication, the new production is
$(110,120)^T+(8/5,6/5)^T=(558/5,606/5)^T$.
Thus the first inverse column records the whole economy's extra
production needed for the one-unit increase in sector 1's final demand.
\end{enumerate}

\exercise{2.6.8}{Prove superposition and interpret an inverse column}
The column-sum hypothesis guarantees $I-C$ invertible for a
nonnegative consumption matrix. For instance if $q<1$ is the
largest column sum, $z=Cz$ would give
$\|z\|_1\le q\|z\|_1$, forcing $z=0$; the square-matrix theorem
then gives invertibility.
\begin{enumerate}[label=(\alph*)]
\item The given equations are $(I-C)x=d$ and
$(I-C)\Delta x=\Delta d$. Adding them gives
$(I-C)(x+\Delta x)=d+\Delta d$. Uniqueness makes
$x+\Delta x$ the new production vector.
\item If $\Delta d=e_1=(1,0,\ldots,0)^T$, then
$\Delta x=(I-C)^{-1}e_1$. Multiplication by $e_1$ selects column
1 of any matrix, so it selects exactly the claimed inverse column.
The same interpretation applies to column $j$ for an increase in
sector $j$'s final demand.
\end{enumerate}

\exercise{2.6.9}{Solve a three-sector numerical production model}
Multiply $(I-C)x=d$ by $10$:
\[8x_1-2x_2=400,\qquad-3x_1+9x_2-3x_3=600,
\qquad-x_1+8x_3=800.\]
The third equation gives $x_1=8x_3-800$. Substituting into the first
gives $x_2=32x_3-3400$. Insert both expressions in the middle equation:
\[-3(8x_3-800)+9(32x_3-3400)-3x_3=600.\]
Thus $261x_3=28800$, giving $x_3=3200/29$,
$x_1=2400/29$, and $x_2=3800/29$.
\answer{$x=(2400/29,3800/29,3200/29)^T$; direct substitution
checks all three equations.}

\exercise{2.6.10}{A demand increase affects every sector when every inverse entry is positive}
An increase $h>0$ in final demand for sector $j$ is $\Delta d=he_j$.
The production change is
$\Delta x=h(I-C)^{-1}e_j$, namely $h$ times inverse column $j$.
Every entry of that column is strictly positive by hypothesis.
\answer{Production must rise in every sector, including sectors whose
final demand did not change directly. Intermediate purchases transmit
the increase throughout this model.}

\exercise{2.6.11}{Verify the two GDP expressions algebraically}
The production equation gives $d=(I-C)x$.
The dual price equation gives $v=p-C^Tp$, so transposing yields
$v^T=p^T-p^TC=p^T(I-C)$. Therefore
\[p^Td=p^T(I-C)x=v^Tx.\]
This uses associativity and the transpose identity
$(C^Tp)^T=p^TC$; all three expressions are scalar values.
It proves the equality without assuming an inverse for either equation.

\exercise{2.6.12}{An iterative inverse from finite geometric sums}
By definition $D_m=I+C+\cdots+C^m$. Multiplying by $C$ shifts
every power up by one, so
\[D_{m+1}=I+CD_m,\qquad D_0=I.\]
Equivalently $D_{m+1}-D_m=C^{m+1}$. Thus repeatedly multiply
the current matrix by $C$ and add $I$ to obtain the next approximation.
To verify the limit, the finite telescoping identity is
$(I-C)D_m=I-C^{m+1}$. The hypothesis $C^m\to0$ implies
$I-C$ invertible: if $z=Cz$, then $z=C^mz\to0$, so $z=0$.
Consequently
\[D_m=(I-C)^{-1}(I-C^{m+1})\longrightarrow(I-C)^{-1}.\]
The recurrence therefore approximates the requested inverse, not just
an unexplained matrix sequence.

\exercise{2.6.13}{Compute production for the supplied seven-sector economy}
\step{Enter the supplied consumption data exactly.}
To specify the data without ambiguous rounding, write
\[C=\frac1{10000}\left[\begin{matrix}1588 & 64 & 25 & 304 & 14 & 83 & 1594\\57 & 2645 & 436 & 99 & 83 & 201 & 3413\\264 & 1506 & 3557 & 139 & 142 & 70 & 236\\3299 & 565 & 495 & 3636 & 204 & 483 & 649\\89 & 81 & 333 & 295 & 3412 & 237 & 20\\1190 & 901 & 996 & 1260 & 1722 & 2368 & 3369\\63 & 126 & 196 & 98 & 64 & 132 & 12\end{matrix}\right].\]
Use the demand vector
$d=(74000,56000,10500,25000,17500,196000,5000)^T$ in the
same sector order as the problem. Solve $Mx=d$ with $M=I-C$;
this is a seven-equation linear system, not a demand for external data.
\step{Solve and check the residual before rounding.}
Exact rational elimination gives the following values, rounded to six
decimal places for display (units: millions of dollars):
\[x\approx\begin{bmatrix}99575.653398\\97703.022863\\51230.523166\\
131569.921929\\49488.491372\\329554.452570\\13835.335715\end{bmatrix}.\]
The exact rational solution satisfies $Mx-d=0$ in every coordinate.
The six-decimal display is only an approximation and need not have
an identically zero numerical residual.
To reproduce that exact calculation, define $C$ using the integer
matrix above divided by $10000$, and run:
\begin{verbatim}
from sympy import Matrix, eye
# C is the displayed integer matrix divided by 10000.
d = Matrix([74000,56000,10500,25000,17500,196000,5000])
M = eye(7) - C
x = M.LUsolve(d)      # exact rational elimination
assert M*x == d
print(x.evalf(12))
\end{verbatim}
This solves the system without constructing an inverse. The displayed
vector is the actual production plan, not merely an algorithm for finding it.

\exercise{2.6.14}{Use the same coefficient matrix for a new demand}
Keep the consumption matrix from Exercise 2.6.13 and change only
the demand to
\[d=(99640,75548,14444,33501,23527,263985,6526)^T.\]
Solving $(I-C)x=d$ by exact arithmetic, then displaying six decimals,
gives (again in millions of dollars)
\[x\approx\begin{bmatrix}134033.615328\\131686.642968\\69471.699079\\
176912.002023\\66595.556706\\443772.902077\\18431.177683\end{bmatrix}.\]
The same reproduction code works with this new $d$.
Checking the exact solution gives $(I-C)x-d=0$ coordinate by coordinate.
A previously computed $LU$ factorization of $I-C$ can be reused;
only its two triangular solves need be repeated.

\exercise{2.6.15}{Count the actual number of iteration steps}
Start with the specified $x^{(0)}=d$ from Exercise 2.6.13 and
iterate $x^{(k+1)}=d+Cx^{(k)}$ without rounding intermediate values.
We interpret ``to four significant figures'' as agreement, in all
seven coordinates, with rounding the exact solution to four significant
figures. The resulting target vector is
\[\begin{bmatrix}9.958\times10^4\\9.770\times10^4\\5.123\times10^4\\
1.316\times10^5\\4.949\times10^4\\3.296\times10^5\\1.384\times10^4\end{bmatrix}.\]
Actual iteration gives these checks:
\[\begin{array}{c|r|c}
k&\max_i |x_i^{(k)}-x_i|/|x_i|&\text{all rounded entries agree?}\\\hline
0&8.09987\times10^{-1}&\text{no}\\
5&4.04043\times10^{-2}&\text{no}\\
10&1.62676\times10^{-3}&\text{no}\\
15&6.49063\times10^{-5}&\text{no}\\
16&3.40686\times10^{-5}&\text{yes}
\end{array}.\]
At step 16, before rounding,
\[x^{(16)}\approx\begin{bmatrix}
99575.105594\\97701.855867\\51228.777814\\131566.362702\\
49487.009311\\329550.480997\\13835.050858\end{bmatrix}.\]
No preceding iterate has all seven required rounded values. At step
15, for instance, the first coordinate is $99574.609463$, which
rounds to $9.957\times10^4$, not $9.958\times10^4$.
Here is a reproducible counting procedure using the exact $C,d$:
\begin{verbatim}
x_exact = (eye(7)-C).LUsolve(d)
four = lambda v: [format(float(z), '.4g') for z in v]
target = four(x_exact)
xk = d
k = 0
while four(xk) != target:
    xk = d + C*xk
    k += 1
assert k == 16
\end{verbatim}
The rational arithmetic prevents iteration-roundoff accumulation;
the final formatting is used only for the specified comparison.
\answer{With $x^{(0)}=d$ counted as step zero, 16 updates are needed.
The usual relative-error check $\max_i |x_i^{(k)}-x_i|/|x_i|<5\times10^{-5}$
also first holds at step 16 for these data.}
